% General MIPLIB-derived instances: recession directions, the HPR bound,
% the HPR -> follower -> tie-break pipeline, coupling rows, tolerances.
\section{General instances: rays, HPR bound and pipelines}\label{app:general}

Notation is that of~\eqref{eq:bl}. In this section a follower row is written
$\ell_i\le C_ix+D_iy\le h_i$ with $\ell_i\in\mathbf{R}\cup\{-\infty\}$ and
$h_i\in\mathbf{R}\cup\{+\infty\}$, and the follower's variable bounds as
$\ell^y_j\le y_j\le h^y_j$, because the direction test below depends on which
sides are finite.

\subsection{Unbounded followers: no bilevel feasible point}

\begin{theorem}[Recession direction]\label{thm:ray}
Let $r\in\mathbf{Q}^{n_y}$ satisfy, for every follower row $i$, $D_ir\ge 0$ if
$\ell_i>-\infty$ and $D_ir\le 0$ if $h_i<+\infty$; for every follower variable
$j$, $r_j\ge 0$ if $\ell^y_j>-\infty$ and $r_j\le 0$ if $h^y_j<+\infty$; and
$f^{\top}r<0$. Then for every leader choice $x$ (real or integral) the
follower's problem $P(x)$ is infeasible or unbounded, so $S(x)=\emptyset$ and
the bilevel problem has no feasible point.
\end{theorem}
\begin{proof}
Let $\nu$ be a positive integer such that $\nu r_j\in\mathbf{Z}$ for every $j\in I_l$
($r$ is rational). Suppose $y$ is feasible for $P(x)$. For every integer
$t\ge 0$, $y+t\nu r$ satisfies every follower row, since $C_ix$ is a constant
and $D_i(y+t\nu r)$ moves only in the permitted direction; it satisfies every
bound for the same reason; and its integer-constrained entries stay integral.
Its objective $f^{\top}y+t\nu f^{\top}r$ tends to $-\infty$. So $P(x)$ has no
optimal solution whenever it is feasible.
\end{proof}
\noindent Not formalised.

\begin{remark}[Meyer's theorem]
The theorem only needs an explicit rational direction. It is the case we
need of a classical result of Meyer:\cite{meyer1974existence} with rational data, the convex hull of the feasible set
of a mixed-integer program is a rational polyhedron with the same recession
cone as the LP relaxation, so a feasible mixed-integer program whose LP
relaxation is unbounded is itself unbounded. Rationality matters: without it,
a direction of the LP relaxation need not be usable by integer points. The
directions here are found by an LP solver in floating point, written as
decimals (which are rationals) and checked exactly.
\end{remark}

\paragraph{The exact check.} For each certificate, a script that reads the
instance files independently of the solver code converts every entry of $r$
to an exact rational and checks the conditions of Theorem~\ref{thm:ray} in
rational arithmetic, counting only follower variables in each follower row
(leader variables enter those rows as constants). When rounding leaves tiny
residuals, the direction is first projected onto the exact null space of its
tight rows. This succeeds for \NumRayExact{} of the \NumNoFeasible{} open instances, whose
certificates are therefore exact (class~(i), Section~\ref{app:classes}); for
\texttt{supportcase42\_0\_100} and \texttt{supportcase42\_50\_50} residuals of
about $10^{-12}$ (relative) remain, so those two certificates rest on floating
point (class~(ii)).
% ROUND2: update when results/rays has exact certificates for supportcase42_0_100 and
% supportcase42_50_50 (exact.ray.json checked by scripts/verify_ray.py): then all
% \NumNoFeasible{} are class (i) and the last clause goes.

\subsection{The HPR bound and the MIPLIB optimum}

\begin{lemma}[HPR bound]\label{lem:hpr}
Every bilevel feasible point is feasible for the HPR, so
$\mathrm{HPR}\le\mathrm{OPT}$. If a bilevel feasible point has value
$\mathrm{HPR}$, it is optimal.
\end{lemma}
\begin{proof}
A bilevel feasible point satisfies $Ax+By\ge a$, the integrality conditions,
and, as $y\in S(x)$, also $Cx+Dy\ge b$. The HPR minimises the same objective
over a superset.
\end{proof}
\noindent Not formalised.

BOBILib builds its general instances from a MIPLIB problem by assigning
variables and rows to the two levels and keeping the problem's objective as
the leader's. If nothing else changes, dropping the condition $y\in S(x)$
gives back the MIPLIB problem: the HPR \emph{is} the MIPLIB model, and its
published optimum\cite{gleixner2021miplib} $z^{*}$ is a lower bound on $\mathrm{OPT}$ by
Lemma~\ref{lem:hpr}. This is not assumed: a script reads the bilevel
instance and the MIPLIB~2017 model and checks in exact rational arithmetic
that they have the same variables, rows, coefficients, sides, bounds,
integrality requirements and objective (up to the objective's sense). For
the five instances where it holds and a verified bilevel feasible point
attains the bound (\texttt{pigeon-10\_0\_100}, \texttt{pigeon-10\_50\_50},
\texttt{timtab1\_0\_100}, \texttt{noswot\_0\_100}, \texttt{noswot\_50\_50}),
the point is optimal by accepted benchmark conventions. The values $z^{*}$
are the optima in MIPLIB~2017's solution file. They were computed by MILP
solvers in floating point and are solver-certified numerical results
(class~(ii), Section~\ref{app:classes}), not exact-arithmetic proofs, so the
five optima are no more rigorous than these published values:
\begin{itemize}
\item \texttt{noswot}: $z^{*}=-41.00000885$. The objective has integer
coefficients on integer variables only, so the HPR optimum is an integer, and
a valid lower bound $L$ gives the lower bound $\lceil L\rceil$. The bound
$\lceil z^{*}\rceil=-41$ therefore holds unless the published optimum
overstates the HPR optimum by $0.99999$ or more (the HPR optimum would have to
be $-42$ or less); BOBILib's own points attain it. Integrality rounding makes
this the safest of the five.
\item \texttt{pigeon-10}: $z^{*}=-9{,}000$, attained by both points; the bound
inherits MIPLIB's optimal status for the instance.
\item \texttt{timtab1}: $z^{*}=764{,}771.99999978$ (printed as $764{,}772$ on
the MIPLIB web pages); the point attains it with row violations below
$10^{-12}$, so the agreement of point and bound is numerical.
\end{itemize}
``Optimal'' for these five therefore means equal to the published optimum to
its printed precision (after rounding up for \texttt{noswot}), as BOBILib and
MIPLIB use the word; no exact-arithmetic proof of the bound is claimed. The
model comparison is exact; the points are class~(ii) (Section~\ref{app:classes}).

\subsection{Feasible points without coupling rows}

\begin{proposition}[HPR, follower, tie-break]\label{prop:general-pipeline}
Assume there are no coupling rows ($B=0$). Let $(\hat x,\hat y)$ be feasible for
the HPR, and suppose $P(\hat x)$ has an optimal solution, with value
$\varphi(\hat x)$. Then the tie-break problem at $\hat x$ (the optimistic choice
among the follower's optima, Section~\ref{app:notation}) is feasible, and for
any of its feasible points $y^{*}$, the point $(\hat x,y^{*})$ is bilevel
feasible. If $y^{*}$ is optimal for the tie-break problem, $(\hat x,y^{*})$ is the
best bilevel feasible point with leader part~$\hat x$.
\end{proposition}
\begin{proof}
$A\hat x\ge a$ holds because $(\hat x,\hat y)$ is HPR feasible and $B=0$, and
$\hat x$ satisfies its integrality conditions. Any optimal solution of
$P(\hat x)$ is feasible for the tie-break problem, whose rows are those of
$P(\hat x)$, $f^{\top}y\le\varphi(\hat x)$, and $A\hat x+By\ge a$, which does not involve
$y$. A feasible point $y^{*}$ of it is feasible for $P(\hat x)$ with
$f^{\top}y^{*}\le\varphi(\hat x)$, so $y^{*}\in S(\hat x)$. The bilevel feasible points
with leader part $\hat x$ are exactly the pairs $(\hat x,y)$ with $y\in S(\hat x)$, the
feasible set of the tie-break problem.
\end{proof}
\noindent Not formalised. The existence of an optimal solution of $P(\hat x)$
is what the follower solve establishes; the HPR point $\hat y$ shows that
$P(\hat x)$ is feasible. The computations accept $f^{\top}y\le\varphi(\hat x)+10^{-7}(1+|\varphi(\hat x)|)$
in the tie-break, well inside BOBILib's relative tolerance~\eqref{eq:bobilib-tol}.
The pipeline handles continuous linking variables, which neither of
BOBILib's two solvers accepts.

\begin{sloppypar}
The tie-break uses this slack. For this supplement the \NumGeneralFiles{} files
of \texttt{results/general} were verified again with \texttt{bobi verify
--generic} (HiGHS as the follower) and the follower value of each file
compared with the follower optimum at the file's leader values. In 19 files
the follower value lies more than $10^{-8}$ (relative) above that optimum:
about $10^{-7}$ in 18 of them (\texttt{binkar10\_1}, \texttt{cmflsp50-24-8-8},
\texttt{core2536-691}, \texttt{cost266-UUE}, \texttt{exp-1-500-5-5},
\texttt{ic97\_potential}, \texttt{n4-3}, \texttt{n5-3}, \texttt{n9-3},
\texttt{neos-2978193-inde}, \texttt{neos-5188808-nattai}, \texttt{neos-873061},
\texttt{neos17}, \texttt{pg}, \texttt{pg5\_34}, \texttt{seymour1},
\texttt{tr12-30}, \texttt{uct-subprob}, all \texttt{\_0\_100}) and one unit,
$5\times10^{-8}$ relative, in \texttt{netdiversion\_0\_100}. These 19 files
are benchmark-accepted points (class~(iii), Section~\ref{app:classes}). The
instances themselves have bilevel feasible points with the same leader values,
since without coupling rows every optimal follower response is feasible for
the leader's rows (Proposition~\ref{prop:general-pipeline}); but the leader
values stated in the 19 files are not shown to be attainable with an exactly
optimal follower. The other files, among them the two optimal points of
Table~\ref{tab:hpr}, have the follower within $10^{-8}$ of its optimum
(class~(ii)).
\end{sloppypar}
% ROUND2: update when results/general records the follower error of each file
% (or the 19 files are re-solved with an exactly optimal follower). The 19 and
% the list come from a re-verification run for this supplement (bobi verify
% --generic, time limit 120 s, on all 64 files), not from the data package.

\subsection{Coupling rows}

With coupling rows ($B\neq 0$) the follower's optima at the HPR's leader
values may all violate them. Recall that $P(x)$, $\varphi(x)$ and $S(x)$ depend
on $x$ only through the linking variables $x_L$; write $\varphi(x_L)$. For fixed
values $\bar x_L$ define the \emph{best response}
\[
  \mathrm{BR}(\bar x_L)=\min\bigl\{c^{\top}x+d^{\top}y:\ x_L=\bar x_L,\ Ax+By\ge a,\
  Cx+Dy\ge b,\ f^{\top}y\le\varphi(\bar x_L),\ \text{integrality}\bigr\},
\]
in which the other leader variables are free.

\begin{proposition}[Best response and blocked leader choices]\label{prop:coupled}
Suppose $P(\bar x_L)$ has an optimal solution. Then the feasible points of
$\mathrm{BR}(\bar x_L)$ are exactly the bilevel feasible points with linking part
$\bar x_L$, and $\mathrm{BR}(\bar x_L)$ is the best value among them. If
$\mathrm{BR}(\bar x_L)$ is infeasible, no bilevel feasible point has linking part
$\bar x_L$ (the choice $\bar x_L$ is \emph{blocked}). Consequently
$\mathrm{OPT}=\inf_{\bar x_L}\mathrm{BR}(\bar x_L)$ over linking values for which
$P(\bar x_L)$ has an optimal solution.
\end{proposition}
\begin{proof}
$(x,y)$ is bilevel feasible iff it satisfies the upper-level rows and
integrality and $y\in S(x)$, and $y\in S(x)$ iff $y$ is feasible for $P(x_L)$ and
$f^{\top}y\le\varphi(x_L)$. These are the constraints of $\mathrm{BR}(x_L)$. If
$P(\bar x_L)$ has no optimal solution, $S(x)=\emptyset$ for every $x$ with these
linking values.
\end{proof}
\noindent Not formalised.

\begin{lemma}[Value-function cuts\cite{mitsos2010global,lozano2017value}]\label{lem:vcut}
Let $\hat y$ be feasible for $P(\hat x_L)$ with $f^{\top}\hat y=\hat f$. Every bilevel
feasible $(x,y)$ satisfies: $\hat y$ is infeasible for $P(x_L)$, or
$f^{\top}y\le\hat f$.
\end{lemma}
\begin{proof}
If $\hat y$ is feasible for $P(x_L)$ then $\varphi(x_L)\le\hat f$, and $y\in S(x)$ gives
$f^{\top}y=\varphi(x_L)$.
\end{proof}
\noindent Not formalised.

The search for an unblocked linking choice tries, in turn within a time
budget: the HPR's own $x_L$; best-response dynamics; the HPR with objective
$c^{\top}x+d^{\top}y+\rho f^{\top}y$ for increasing $\rho$; a penalty alternating
direction method after Kleinert and Schmidt;\cite{kleinert2021computing} moves of continuous linking
variables towards the interior; the HPR with the cuts of
Lemma~\ref{lem:vcut} (modelled with binary indicators, exact when every
follower row's linking part takes integer values); and no-good cuts on
integer linking variables. These are heuristics: none of them is needed for
correctness, and the ``blocked'' verdicts (solver tolerances included) only
steer the search; no instance is claimed to have no feasible point on that
basis. What is certified is each returned point, re-checked from scratch: all
rows, bounds and integrality within $10^{-6}$; a fresh follower solve whose
proven dual bound puts the point's follower value within relative $10^{-5}$ of
the optimum; and the optimistic tie-break at the point's leader values within
relative $10^{-5}$ by a proven bound. Each file then passes \texttt{bobi
verify}.

% ROUND2: update when the 9 coupled files in results/general-coupled are repaired or
% re-classified (they are class (iii) until then).
In \NumCoupledTolerance{} of the \NumCoupledFiles{} files with coupling rows the follower's value lies about
$10^{-7}$ (relative) above its optimum, which BOBILib's
test~\eqref{eq:bobilib-tol} accepts; these files are benchmark-accepted points
(class~(iii), Section~\ref{app:classes}), and the other files with coupling
rows are class~(ii). For seven of them a point with the
follower at its optimum (within $10^{-9}$ relative) exists at the same leader
values with a slightly worse leader objective. For \texttt{pg\_50\_50} and
\texttt{neos-2978193-inde\_50\_50} the coupling rows force the follower about
$10^{-7}$ above its optimum at the files' leader values; with the non-linking
leader variables free, exactly optimal responses exist, so both instances do
have bilevel feasible points, but the two files are feasible only within the
tolerance.

\paragraph{Few binary linking variables.} If the linking variables are
binary and few, Proposition~\ref{prop:coupled} gives the exact optimum by
enumerating all $\bar x_L$, solving $P(\bar x_L)$ and $\mathrm{BR}(\bar x_L)$ for
each. This reproduces the published optimum $299{,}577$ of Lozano, Bergman
and Cire\cite{lozano2026network} for \texttt{general30-30-10-20-20-8} (10
linking bits); with 20 bits it is too slow.

\subsection{A tolerance observation}

\begin{remark}[\texttt{general20-30-20-20-20-2}]\label{rem:tolerance}
Test~\eqref{eq:bobilib-tol} is relative. In the generalized interdiction
instance \texttt{general20-30-20-20-20-2} the follower's objective values are
near $10^8$, so a response up to about $10^4$ units from the follower's optimum
passes. In BOBILib's solution file for this instance the follower's value is
$106{,}899{,}984.14$ against the follower's optimum $106{,}899{,}953$ at the same
leader choice, $31$ units above it (relative error about $2.9\times10^{-7}$), and
rows are violated by up to $3.3\times10^{-4}$ (our verifier's record is in
\texttt{results/notes/}). The row violations pass the $10^{-4}$ tolerance of
the MIPLIB checker inside BOBILib's checker (Section~\ref{app:tolerances}) but
not the $10^{-6}$ of \texttt{bobi verify}, and the stated objective
\GenTolObjFile{} differs from the \GenTolObjComputed{} of the point, so
\texttt{bobi verify} rejects the file. Such a point is
not bilevel feasible in the sense of Definition~\ref{def:bilevel-feasible}:
with an exactly optimal follower response the leader's value at the same
leader choice can differ. A value accepted only under a relative tolerance
is therefore not proven attainable with an exactly optimal follower
response: the file is the example of a benchmark-accepted point
(class~(iii), Section~\ref{app:classes}).
\end{remark}
